\section{\esp INTRODUÇÃO}

O marketing digital é uma atividade estratégica para organizações que desejam divulgar produtos e serviços, alcançar públicos específicos e acompanhar resultados por indicadores mensuráveis \cite{gujar2024,younas2025}. No âmbito da publicidade digital paga, a cadeia de serviços envolve anunciantes, agências de mídia, \,\textit{publishers} digitais e plataformas de troca de anúncios, entre os quais circulam informação, recursos financeiros e serviços \cite{agus2019}.

Em operações de marketing digital, a gestão de campanhas envolve criação, orçamento, públicos, conteúdo e análise de resultados \cite{younas2025,yahia2024,nain2025}. É nesse contexto operacional que se insere a proposta do SIGE Desk, uma aplicação web de Help Desk destinada a organizar demandas de tráfego pago. No sistema, esses elementos poderão ser tratados como solicitações, tais como criação ou pausa de anúncios, ajustes de orçamento e de público, aprovação de criativos e pedido de relatórios. A literatura de Service Desk relata que o encaminhamento de solicitações diretamente a diferentes pessoas pode produzir múltiplos pontos de contato e apresenta o Service Desk como ponto único para receber, registrar, analisar, encaminhar e acompanhar solicitações \cite{firmansyah2022,fenner2015,amanullah2017}.

Os estudos revisados sobre marketing digital concentram-se na automação, no planejamento, no orçamento e no desempenho das campanhas \cite{gujar2024,younas2025,yahia2024,nain2025}. Os estudos de Service Desk, por sua vez, tratam do recebimento e do acompanhamento de solicitações \cite{firmansyah2022,fenner2015,amanullah2017}. A proposta do SIGE Desk aproxima esses campos ao tratar uma demanda de tráfego pago como solicitação de serviço a ser registrada e acompanhada, sem executar ou automatizar a campanha.

Quando as informações de uma solicitação não são mantidas em um registro comum, tornam-se mais difíceis a identificação do pedido original, a definição de responsável, o acompanhamento de prazo, o registro de aprovação e a recuperação das decisões tomadas \cite{fenner2015,servicedesk2023,servicerequest2023}. Essa dificuldade é particularmente relevante no tráfego pago, pois alterações de orçamento, público ou criativo podem afetar o desempenho e o custo de uma campanha \cite{yahia2024}. As práticas de Service Desk e de gerenciamento de solicitações procuram mitigar esses riscos por meio de registro, triagem, comunicação e acompanhamento até a conclusão \cite{firmansyah2022,servicedesk2023,servicerequest2023}. Assim, a questão que orienta o trabalho é: \textit{como uma aplicação web de Help Desk pode centralizar e rastrear demandas de tráfego pago, melhorando o acompanhamento de campanhas e o atendimento a clientes de agências de marketing digital?}

O estudo justifica-se pela possibilidade de organizar as solicitações em tickets com responsáveis, prioridades, prazos, comentários, aprovações, histórico e evidências. Essa estrutura pode apoiar a distribuição do trabalho, a prestação de contas ao cliente e a consulta a indicadores de atendimento, sem substituir as plataformas de mídia nem automatizar a alteração dos anúncios. Do ponto de vista acadêmico, a proposta integra Sistemas de Informação, Gestão de Serviços, Gestão de Projetos e Marketing Digital, aplicando conhecimentos de levantamento de requisitos, modelagem de processos, experiência do usuário, desenvolvimento web e validação de software a um problema de gestão.

O objetivo geral é desenvolver uma aplicação web de Help Desk para centralizar, acompanhar e registrar demandas de tráfego pago. Para tanto, pretende-se revisar artigos e normas aplicáveis; levantar, analisar e priorizar requisitos; modelar o processo de atendimento e o backlog; prototipar e implementar um produto mínimo viável com tickets, histórico, prazos e comunicação entre os envolvidos; e realizar testes funcionais e avaliação de usabilidade.
